\section{随机化 SVD：用随机向量探到主子空间}
\label{sec:08-Numerical-Analysis/05-Numerical-Linear-Algebra/09}

一个 \(10^6\times10^5\) 的嵌入矩阵，要取前 \(k=200\) 个奇异方向做降维。完整 SVD 的运算量是 \(O(mn^2)\approx10^{16}\)，机器跑不动，而且你根本不需要后面那九万多个方向。退一步用\hyperref[sec:08-Numerical-Analysis/05-Numerical-Linear-Algebra/04]{``大型稀疏系统与 Krylov 迭代''}那一路的 Lanczos 双对角化，只用矩阵--向量乘，理论上合适；实际跑起来收敛很慢，因为这个矩阵的谱在前几百个位置上挤得很密，Krylov 子空间要长到很深才分得开，而每深一层都要对已有的基重新正交化。

问题不在算力，在于这两条路都想\emph{精确地}把前 \(k\) 个方向一个一个揪出来。可你要的其实不是这些方向本身，是它们张成的那个 \(k\) 维子空间。找子空间比找方向容易得多，而找子空间恰好是随机向量擅长的事。

\subsection{三步：探、正交化、还原}

记 \(\ell=k+p\)，\(p\) 是一个小的过采样量。

第一步，取 \(\Omega\in\R^{n\times\ell}\)，元素独立取自 \(N(0,1)\)，算 \(Y=A\Omega\)。\(Y\) 的每一列是 \(A\) 作用在一个随机向量上的结果，因而落在 \(A\) 的列空间里，而且——下一小节要说明——以高概率偏向大奇异值那一侧。

第二步，对 \(Y\) 做 QR 分解 \(Y=QR\)，扔掉 \(R\)。\(Q\in\R^{m\times\ell}\) 的列是 \(Y\) 列空间的一组标准正交基，也就是我们探到的那个近似主子空间。

第三步，令 \(B=Q^{\trans}A\in\R^{\ell\times n}\)，对这个\emph{小}矩阵做精确 SVD，\(B=\tilde U\Sigma V^{\trans}\)，再转回去 \(U=Q\tilde U\)。于是
\[
A\;\approx\;QQ^{\trans}A\;=\;Q B\;=\;U\Sigma V^{\trans},
\]
截取前 \(k\) 列就是要的东西。全程只对 \(A\) 做了两次遍历（算 \(Y\)、算 \(B\)），此外全是 \(O(m\ell^2+n\ell^2)\) 的小规模运算。

\subsection{随机向量为什么能探到主子空间}

这是全篇的直觉支点，值得说透。把随机向量 \(\omega\)（\(\Omega\) 的某一列）按右奇异向量展开，\(\omega=\sum_j c_jv_j\)。\(\Omega\) 的元素独立同分布且各向同性，于是系数 \(c_j=v_j^{\trans}\omega\) 全部服从同一个 \(N(0,1)\)——\emph{随机向量对所有方向一视同仁，它不偏爱任何一个 \(v_j\)}。

现在让 \(A\) 作用上去：
\[
A\omega=\sum_j\sigma_jc_ju_j .
\]
在 \(u_j\) 方向上的分量大小是 \(\sigma_j\abs{c_j}\)。既然 \(\abs{c_j}\) 同分布，两个方向的分量之比在典型情形下就是 \(\sigma_j/\sigma_{j'}\)。一次矩阵乘法之后，原本各向同性的向量已经被拉成了偏向大奇异值的形状：\(\sigma_1\) 是 \(\sigma_{100}\) 的十倍，那第一个方向上的分量典型地也就是第一百个方向上的十倍。

不需要哪一列特别``好''。要的只是 \(\ell\) 列合起来张成的空间盖住前 \(k\) 个方向，而每一列都在往同一侧倾斜，多几列就把这件事做扎实了——过采样的 \(p\) 列买的正是这份余量。反过来，若 \(\sigma_k\) 与 \(\sigma_{k+1}\) 几乎一样大，倾斜就不明显，这个机制立刻变弱。误差界里出现的正是这件事。

\subsection{误差界的形式与读法}

\begin{proposition}[随机化取子空间的期望误差]
设 \(\Omega\) 元素独立取自 \(N(0,1)\)，\(\ell=k+p\) 且 \(p\ge2\)，\(Q\) 由 \(Y=A\Omega\) 的 QR 分解得到，则
\[
\E\norm{A-QQ^{\trans}A}_2
\;\lesssim\;\Bigl(1+\sqrt{k/p}\Bigr)\,\sigma_{k+1}
\;+\;\frac{e\sqrt{k+p}}{p}\Bigl(\sum_{j>k}\sigma_j^2\Bigr)^{1/2}.
\]
\end{proposition}

不证，只读。第一项是主项：\(\sigma_{k+1}\) 是任何秩 \(k\) 逼近都逃不掉的下界（Eckart--Young，见\hyperref[sec:08-Numerical-Analysis/05-Numerical-Linear-Algebra/06]{``SVD、低秩与数值秩''}），随机化在它前面乘了一个 \(1+\sqrt{k/p}\)。\(k=200\)、\(p=10\) 时这个因子约为 5.5，\(p=20\) 时约为 4.2——注意 \(p\) 是加在 \(\ell\) 上的常数，从 5 加到 20 几乎不改变运算量，却把常数明显压下来。这就是``\(p\) 取 5 到 10''这条经验的来源：它不是调参玄学，是一个可以事先算出来的、随 \(p\) 平方根衰减的因子。这里的取法有理论依据，具体取 5 还是 10 则是在多次实践中稳定下来的默认值，属于经验规律，换一类矩阵不保证最优。

第二项是尾部项，它把 \(\sigma_{k+1}\) 之后所有奇异值的平方和都算了进来。谱衰减快时这一项可以忽略，衰减慢时它会盖过主项。

\subsection{谱衰减慢，界就塌了}

\begin{center}
\begin{tikzpicture}[x=1cm,y=1cm,font=\small,>=stealth]
  \draw[black!55,->] (0,0) -- (5.0,0);
  \draw[black!55,->] (0,0) -- (0,3.5);
  \draw[thick,blue!62] plot[smooth] coordinates
    {(0.35,3.00) (0.73,1.86) (1.11,1.15) (1.49,0.72) (1.87,0.44) (2.25,0.28)
     (2.63,0.17) (3.01,0.11) (3.39,0.07) (3.77,0.04) (4.15,0.03) (4.53,0.02)};
  \draw[thick,orange!85!black] plot[smooth] coordinates
    {(0.35,3.00) (0.73,2.20) (1.11,1.85) (1.49,1.61) (1.87,1.44) (2.25,1.32)
     (2.63,1.22) (3.01,1.14) (3.39,1.07) (3.77,1.02) (4.15,0.97) (4.53,0.92)};
  \draw[dashed,black!55] (1.87,0) -- (1.87,1.70);
  \node[below,font=\scriptsize,black!70] at (1.87,-0.06) {\(k+1\)};
  \draw[<->,red!70!black,thick] (2.05,0.44) -- (2.05,1.44);
  \node[right,align=left,font=\scriptsize,red!70!black] at (2.15,0.80)
    {同一个 \(k\)，\\两种谱下的 \(\sigma_{k+1}\)};
  \node[right,font=\scriptsize,blue!62] at (3.45,0.38) {衰减快};
  \node[right,font=\scriptsize,orange!85!black] at (3.45,1.55) {衰减慢};
  \node[below,font=\footnotesize,black!70] at (2.5,-0.55) {奇异值序号 \(i\)};
  \node[left,font=\scriptsize,black!65] at (-0.05,3.2) {\(\sigma_i\)};
  \begin{scope}[shift={(6.3,0)}]
    \draw[black!55,->] (0,0) -- (5.0,0);
    \draw[black!55,->] (0,0) -- (0,3.5);
    \draw[dashed,thick,black!35] plot[smooth] coordinates
      {(0.35,3.00) (0.73,2.20) (1.11,1.85) (1.49,1.61) (1.87,1.44) (2.25,1.32)
       (2.63,1.22) (3.01,1.14) (3.39,1.07) (3.77,1.02) (4.15,0.97) (4.53,0.92)};
    \draw[thick,green!45!black] plot[smooth] coordinates
      {(0.35,3.00) (0.73,1.18) (1.11,0.69) (1.49,0.46) (1.87,0.34) (2.25,0.26)
       (2.63,0.21) (3.01,0.18) (3.39,0.15) (3.77,0.13) (4.15,0.11) (4.53,0.10)};
    \node[right,font=\scriptsize,black!45] at (3.40,1.40) {原谱};
    \node[right,align=left,font=\scriptsize,green!40!black] at (2.30,0.62)
      {幂迭代之后：\\\(\sigma_i^{2q+1}\)（归一化）};
    \node[below,font=\footnotesize,black!70] at (2.5,-0.55) {奇异值序号 \(i\)};
  \end{scope}
\end{tikzpicture}

\emph{左：同一个截断位置 \(k\)，衰减快的谱把 \(\sigma_{k+1}\) 压到很低，误差界的主项跟着变小；衰减慢的谱在同一位置还挂在高处，那条界即使常数为 1 也不好看——问题出在矩阵身上，不在算法身上。右：幂迭代把每个奇异值取 \(2q+1\) 次方并重新归一化，慢衰减的曲线被人为拉陡，橙线那种情形于是被搬到蓝线一侧。}
\end{center}

界里的 \(\sigma_{k+1}\) 是矩阵自己的性质，算法改不了。能改的是矩阵——把 \(A\) 换成 \((AA^{\trans})^qA\)，它的奇异向量与 \(A\) 完全相同，奇异值变成 \(\sigma_i^{2q+1}\)。于是
\[
Y=(AA^{\trans})^qA\Omega ,
\]
相对间隔从 \(\sigma_{k+1}/\sigma_k\) 变成 \((\sigma_{k+1}/\sigma_k)^{2q+1}\)。\(\sigma_{k+1}/\sigma_k=0.9\) 这种几乎分不开的情形，\(q=3\) 之后变成 \(0.9^{7}\approx0.48\)，一下子拉开了。误差界里那个 \(1+\sqrt{k/p}\) 的因子相应地开 \(2q+1\) 次方，\(q=2\) 就足以把它压到 1.4 以下。

这个动作与\hyperref[sec:08-Numerical-Analysis/05-Numerical-Linear-Algebra/05]{``从幂迭代到 QR 特征值算法''}里的幂迭代是同一回事，区别只在这里同时推着 \(\ell\) 个向量走，而不是一个。

\subsection{条件对账}

\textbf{要付出的是 \(2q\) 次额外的矩阵乘法。}\(q=2\) 意味着对 \(A\) 多遍历四次。矩阵在内存里时这不算什么，在磁盘或分布式存储上时每一次遍历都是实打实的开销，\(q\) 就得省着用。所以先看谱：谱衰减快就别加幂迭代，加了纯属浪费。

\textbf{每一步必须重新正交化，否则拉开的间隔会被浮点吞掉。}把 \((AA^{\trans})^q A\Omega\) 按字面顺序连乘是错的实现。\(\ell\) 列在乘的过程中一起向 \(u_1\) 倾斜，到 \(q=2\) 时第 \(j\) 列与第一列的夹角已经小到在双精度下分辨不出——所有列在数值上塌成同一个方向，\(Y\) 的秩掉到 1，QR 分解出来的后 \(\ell-1\) 列全是舍入噪声。正确写法是每乘一次就做一次 QR：
\[
Y_0=A\Omega,\quad Q_0=\operatorname{qr}(Y_0),\quad
Q_i'=\operatorname{qr}(A^{\trans}Q_{i-1}),\quad Q_i=\operatorname{qr}(AQ_i') .
\]
量化一下失效点：\(\sigma_j/\sigma_1=10^{-3}\) 且 \(q=3\) 时，\(\sigma_j^{7}/\sigma_1^{7}=10^{-21}\)，已经在双精度相对精度 \(\varepsilon\approx10^{-16}\) 以下——不重新正交化，这个方向就不再存在于数据里了。\hyperref[sec:08-Numerical-Analysis/01-Floating-Point-and-Error/02]{``误差如何产生与放大''}讲的是同一种失效：有效信息被同一个数量级上的舍入误差淹没，而且没有任何报错。

\textbf{界是对随机性取期望，不是对每次运行成立。}上面那条命题给的是 \(\E\norm{A-QQ^{\trans}A}_2\)。单次运行可以比它差，虽然同类分析也给得出指数型的尾部界（超出期望多少的概率随 \(p\) 指数下降）。要的是可复现的数字就固定随机种子，并且把种子写进报告。

\textbf{保住的是子空间，不是单个奇异向量。}命题控制的是投影残差 \(\norm{A-QQ^{\trans}A}_2\)。\(\sigma_k\) 与 \(\sigma_{k+1}\) 接近时，前 \(k\) 个奇异向量本身在对应的子空间里可以大幅旋转，这一点在\hyperref[sec:08-Numerical-Analysis/05-Numerical-Linear-Algebra/06]{``SVD、低秩与数值秩''}里已经说过，随机化不会让它变好也不会让它变坏。要逐个解释奇异向量的场合，先检查那里有没有谱间隙。

\textbf{输出的形状是 SVD，身份不是 SVD。}拿到的是 \(QQ^{\trans}A\) 的精确分解，不是 \(A\) 的。报告时要附上 \(\norm{A-QQ^{\trans}A}\) 的估计（用几个新的随机向量测残差即可，这本身又是一个随机化估计）、\(\ell\)、\(q\) 和种子。不写这些，下游没有办法判断这份分解能信到什么程度。

\textbf{矩阵不大时别用。}\(m,n\) 只有几千，LAPACK 的稠密 SVD 几秒就出结果，还给的是确定性的、逐个奇异值都有 Weyl 界兜底的答案。随机化的常数开销与额外的验证工作在这个规模上都不划算。

\subsection{一个左乘，一个右乘}

把这一篇和\hyperref[sec:08-Numerical-Analysis/05-Numerical-Linear-Algebra/08]{``用随机草图把大矩阵压小''}并排看，两者用的是同一个随机矩阵，站的位置却相反。草图从左边乘，把行数从 \(m\) 压到 \(s\)，保住的是列空间上的长度，为的是让最小二乘的方程变少。随机化 SVD 从右边乘，把列数从 \(n\) 压到 \(\ell\)，保住的是前 \(k\) 个左奇异方向张成的子空间，为的是让待分解的矩阵变窄。一个压方程，一个探方向。

共同点比这个差别更重要：两者给的都是``以高概率''而不是``一定''。这个词换掉之后，一条结论还能被怎样使用、不能被怎样使用，是下一篇要单独处理的事。
